%% Portfolio Index - Marcus D. Yang
%% Multi-page project summary to attach where an application requires a portfolio
%% file (PDF) but the real portfolio is a website. Gives a reviewer something
%% concrete to look at and drives clicks to the full site. No longer a strict
%% 1-page document - content is allowed to flow to 2-3 pages naturally.
%%
%% Style matches cv/main_spacex_falcon9_intern.tex (one-page-classic fonts/rules),
%% just without the 1-page constraint.
%% Compile with: lualatex -interaction=nonstopmode portfolio_index.tex
%%
%% IMAGES in portfolio/img/ (all present as of this revision):
%%   tvc.jpg, tvc2.jpg, tvc3.jpg, tvc4.jpg
%%     - TVC Rocket: assembled gimbal hardware, wired avionics (IMU/Teensy),
%%       3D printer mid-print (nose cone), and motor-loading/field-test prep.
%%       Pulled from projects-website-psi.vercel.app/projects/tvc-rocket.
%%   f101.jpg
%%     - F-101 Voodoo CAD render, side profile (from FinalProjectAssemblyRight.tif)
%%   liquidrocketry1.jpg, liquidrocketry2.jpg
%%     - Liquid Rocketry at Illinois: P&ID / pressurant-tank assembly CAD
%%       screenshots (Siemens NX), user-supplied
%% Until a named file exists, a labeled placeholder box renders in its place,
%% so this still compiles cleanly if an image goes missing.
%%
%% PORTFOLIO URL: \portfoliourl below.

\documentclass[11pt,letterpaper]{article}

\usepackage{fontspec}
\setmainfont{TeX Gyre Termes}
\usepackage[top=0.6in, bottom=0.6in, left=0.75in, right=0.75in]{geometry}
\usepackage{titlesec}
\usepackage{enumitem}
\usepackage{graphicx}
\usepackage{needspace}
\usepackage[hidelinks]{hyperref}

\setlength{\parindent}{0pt}

% ---- real portfolio URL ----
\newcommand{\portfoliourl}{https://projects-website-psi.vercel.app/}
\newcommand{\portfoliourldisplay}{projects-website-psi.vercel.app}
% ----------------------------

\titleformat{\section}{\raggedright\large\bfseries}{}{0em}{}[\titlerule\vspace{-4pt}]
\titlespacing*{\section}{0pt}{12pt}{6pt}
\setlist[itemize]{leftmargin=0.18in, itemsep=2pt, parsep=0pt, topsep=3pt, partopsep=0pt}

% Image box: #1 basename, #2 width (fraction of \linewidth, as a decimal), #3 max height
\newcommand{\imgbox}[3]{%
  \begin{minipage}[t]{#2\linewidth}
    \centering
    \IfFileExists{img/#1.jpg}
      {\includegraphics[width=\linewidth, height=#3, keepaspectratio]{img/#1.jpg}}
      {\fbox{\parbox[c][0.9\dimexpr#3\relax][c]{0.85\linewidth}{\centering\footnotesize\itshape add img/#1.jpg}}}
  \end{minipage}
}

% Project text block: #1 title, #2 subtitle/date line, #3 bullets
\newcommand{\projecttext}[3]{%
  \textbf{#1}\\[1pt]
  \textit{\footnotesize #2}
  \begin{itemize}
    #3
  \end{itemize}
}

\begin{document}

\begin{center}
  {\Large\textbf{Marcus D. Yang}} \\[3pt]
  Project Portfolio \ $|$ \ \href{\portfoliourl}{\portfoliourldisplay} \ $|$ \ \href{https://github.com/mdyang26}{github.com/mdyang26} \\[2pt]
  {\footnotesize (650) 954-0288 \ $|$ \ \href{mailto:yangdmarcus@gmail.com}{yangdmarcus@gmail.com} \ $|$ \ \href{https://www.linkedin.com/in/marcusdyang}{linkedin.com/in/marcusdyang}}
\end{center}
\vspace{2pt}

{\small Aerospace Engineering student at UIUC. Selected hands-on design-build-test projects below; full write-ups, code, and additional work at \href{\portfoliourl}{\portfoliourldisplay}.}

\section{Selected Projects}

% ============================================================
%     TVC ROCKET (flagship - 4 photos)
% ============================================================
\projecttext{TVC Rocket -- Solo Design, Build \& Test}{Siemens NX, C++ (Teensy 4.1), Python, FDM 3D printing \ $\cdot$ \ Summer 2026}{
  \item Two-axis gimbaled thrust-vector-control rocket, built end to end solo: CAD, embedded PID firmware, 3D-printed structures, static-fire testing.
  \item Monte Carlo sensitivity analysis (hundreds of runs) identified IMU accelerometer bias, not servo speed, as the dominant attitude-error driver.
  \item Root-caused two static-fire failures purely from telemetry, video, and audio data: an estimator axis-convention bug and an inverted control sign.
}
\vspace{4pt}

\noindent
\imgbox{tvc3}{0.235}{1.7in}\hfill
\imgbox{tvc2}{0.235}{1.7in}\hfill
\imgbox{tvc}{0.235}{1.7in}\hfill
\imgbox{tvc4}{0.235}{1.7in}

\vspace{2pt}
{\footnotesize\itshape Left to right: printing the nose cone, wired avionics (IMU + Teensy 4.1), assembled gimbal, loading the motor for a static-fire test.}

\vspace{16pt}

% ============================================================
%     F-101 VOODOO
% ============================================================
\needspace{7\baselineskip}
\projecttext{F-101 Voodoo -- CAD Reconstruction}{Siemens NX \ $\cdot$ \ Spring 2025}{
  \item Rebuilt a full McDonnell F-101 Voodoo from multi-view engineering drawings, scaling dimensions to reconstruct fuselage, wing, and tail geometry.
  \item Integrated articulated control surfaces (ailerons, elevators, rudder) with parametric assemblies and constraints.
}
\vspace{4pt}

\begin{center}
  \imgbox{f101}{0.6}{2.2in}
\end{center}

\vspace{16pt}

% ============================================================
%     LIQUID ROCKETRY AT ILLINOIS
% ============================================================
\needspace{8\baselineskip}
\projecttext{Liquid Rocketry at Illinois -- Member}{Siemens NX, Ansys Fluent, KiCad \ $\cdot$ \ Spring 2025 -- Present}{
  \item Modeled P\&ID system architecture in Siemens NX and integrated components into the pressurant tank assembly.
  \item Created CAD geometries for canard-grid fin configurations and ran Ansys Fluent CFD cases across a 0-20 degree deflection range.
  \item Engineered a KiCad PCB radio transmitter (USB-C + MCU) prototype for ground-test telemetry links.
}
\vspace{4pt}

\noindent
\imgbox{liquidrocketry1}{0.47}{2.6in}\hfill
\imgbox{liquidrocketry2}{0.47}{2.6in}

\vspace{2pt}
{\footnotesize\itshape Pressurant tank / P\&ID assembly modeled in Siemens NX.}

\vspace{20pt}
\begin{center}
  {\small Full portfolio: \ \href{\portfoliourl}{\textbf{\portfoliourldisplay}}}
\end{center}

\end{document}
